\documentclass[../main.tex]{subfiles}
\begin{document}

\begin{center}
    \Large{\textbf{ABSTRACT}}\\
\end{center}
\vspace{1cm}

Trong bối cảnh công nghiệp hiện đại, sự hội tụ giữa mạng công nghệ thông tin (IT) và công nghệ vận hành (OT) trước yêu cầu bắt buộc của quy trình quản lý sản xuất quy mô lớn đã đồng thời làm gia tăng các mối đe dọa an ninh mạng nhắm vào hệ thống điều khiển công nghiệp (ICS/OT). Các cuộc tấn công vào OT gây ra các tác động trong thế giới vật lý, từ nghiêm trọng thấp như gây dừng dây chuyền dẫn tới thiệt hại kinh tế tới nghiêm trọng cao nhất là nguy hiểm đến tính mạng con người đang vận hành dây chuyền. Do đặc thù các thiết bị OT vốn được thiết kế từ đầu với một triết lý duy nhất là ưu tiên tính ổn định, thời gian thực và tuổi thọ cao để có thể hoạt động 24/7 trong suốt vòng đời hoạt động 20-30 năm của nhà máy, nhiều hệ thống hiện nay như các dòng PLC Siemens S7-300/400 vẫn sử dụng giao thức S7comm lạc hậu không tích hợp sẵn cơ chế xác thực hay mã hóa. Trong khi đó, các giải pháp bảo mật truyền thống từ góc nhìn IT để chống lại các cuộc tấn công vào giao thức mạng IT tỏ ra kém hiệu quả trước các giao thức công nghiệp đặc thù và đóng kín của OT.

Vì vậy, đồ án lựa chọn hướng tiếp cận phát hiện xâm nhập thụ động dựa trên lưu lượng mạng của giao thức S7 Comm. Cách tiếp cận này phù hợp với môi trường OT vì không can thiệp trực tiếp vào quá trình điều khiển, đồng thời vẫn cho phép quan sát và phân tích các hành vi bất thường trong giao tiếp giữa các thành phần hệ thống, đồng thời được tinh chỉnh để hoạt động tối ưu trên giao thức công nghiệp đặc thú Siemens S7. Trên cơ sở đó, đồ án triển khai môi trường mô phỏng mạng công nghiệp và các kịch bản tấn công điển hình trên S7comm: trinh sát, ghi biến quá trình (Data Block), từ chối dịch vụ (flood TCP/102 và lệnh S7 hợp lệ), phát lại Start/Stop và quan sát man-in-the-middle. Các kịch bản DoS được gói trong một script (\texttt{--mode all}, có khoảng nghỉ giữa các phase) để triển khai lab và kích hoạt luật tần suất trên IDS. Đồ án xây dựng hai tập luật Suricata---\texttt{s7comm.rules} (opcode/payload) và \texttt{ics\_dos.rules} (\texttt{threshold})---dựa trên mã chức năng, offset byte và mật độ gói tới PLC.

Đóng góp chính là hệ thống hóa kỹ thuật tấn công S7comm trong lab và bộ luật phát hiện tường minh, có thể kiểm chứng bằng Suricata trên lưu lượng mirror. Kết quả cho thấy IDS có thể cảnh báo các hành vi nguy hiểm đã mô phỏng, gồm flood DoS, với tác động tối thiểu tới vận hành sản xuất khi triển khai thụ động.

\begin{flushright}
\begin{tabular}{@{}c@{}}
Student\\
\textit{(Signature and full name)}
\end{tabular}
\end{flushright}

\end{document}